\section*{Pad Thai}

\ihead{}\chead{Personen:2}\ohead{}
\ifoot{Vorbereitungszeit:15 min}\cfoot{}\ofoot{Kochzeit:10 min}

{\Large Zutaten}

\begin{multicols}{2}
\begin{itemize}
    \item \SI{100}{g} Reisnudeln (3 mm breit)
    \item \SI{100}{g} Tofu
    \item \SI{100}{g} Hühnerfleisch
    \item \num{4} Knoblauchzehen
    \item \num{2} Frühlingszwiebeln
    \item \SI{100}{g} Sprossen
    \item \num{2} Eier
    \item \SI{2}{EL} neutrales Pflanzenöl
    \item \SI{2}{TL} brauner Zucker
    \item \SI{4}{TL} Fischsauce
    \item \SI{6}{TL} Austernsauce
    \item \SI{1}{Tasse} Wasser
    \item etwas ungesalzene, geröstete Erdnüsse
    \item etwas Chiliflocken
    \item \num{1} Limette
    \item etwas Kurkumapulver
\end{itemize}
% \columnbreak
\end{multicols}

\noindent {\Large Zubereitung}\\ \\

\textit{Vorbereitung:} Trockene Reisnudeln in heißem Wasser einige Minuten einweichen. 
Tofu mit Kurkuma einreiben und in Stücke schneiden. 
Hähnchenfleisch sehr klein schneiden, Knoblauch fein hacken und Frühlingszwiebeln in etwa \SI{2}{cm} lange Stücke schneiden. 
Sprossen waschen und abtropfen lassen. Erdnüsse hacken und Limette vierteln.\\

Öl im Wok erhitzen und Knoblauch kurz anbraten. 
Hähnchen und kurz darauf Tofu zugeben, etwa \SI{1}{min} braten und zur Seite schieben. 
Eier zugeben, Fischsauce, Zucker und Austernsauce darübergeben und kräftig verrühren. 
Alles vermischen, etwas Wasser und die Nudeln zugeben. 
Sprossen und Frühlingszwiebeln hinzufügen und einige Sekunden mitkochen. 
Sofort mit gehackten Erdnüssen, Chiliflocken und einem Limettenviertel servieren.
